\heading{Seminar 1}{Preferences, utility and counterexamples}
This worksheet accompanies Lecture 1. No demand calculus is needed. Give definitions before calculations, and justify every answer. For a false universal claim, one fully specified counterexample is enough. The final exercise is reserve material.
\conventions
\task{1. Relations before utility}{Core · 15 minutes}
On $X=\R_+^2$, classify each relation as complete or incomplete, transitive or intransitive, continuous or discontinuous. Prove every positive claim; give a numerical witness for every failure.
\begin{enumerate}[label=(\alph*)]
\item $x\pref_A y$ iff $x_1\ge y_1$ and $x_2\ge y_2$.
\item $x\pref_B y$ iff $x_1\ge y_1$ or $x_2\ge y_2$.
\item $x\pref_C y$ iff $x_1+x_2>y_1+y_2$, or both sums are equal and $x_1\ge y_1$.
\end{enumerate}
Is relation C just an ordinary linear utility with a sufficiently small positive weight added to good 1?
\task{2. True or false?}{Core · 15 minutes}
\begin{enumerate}[label=(\alph*)]
\item If two bundles are incomparable, the consumer is indifferent between them.
\item If a relation is represented by a real-valued function, it must be rational.
\item If a representing utility is discontinuous, the preferences must be discontinuous.
\item A convex utility function represents convex preferences.
\item Strictly increasing transformations preserve convexity of preferences, but need not preserve concavity of utility.
\item LNS implies that utility is unbounded above.
\end{enumerate}
\task{3. A kink, a transformation and a curve}{Core · 20 minutes}
Let $X=\R_+^2$, $u(x)=\min\{2x_1,x_2\}$ and $v(x)=e^{u(x)}$.
\begin{enumerate}[label=(\alph*)]
\item Draw the indifference set through $(1,2)$ and its weak upper contour set. Include the corner and both rays.
\item Are the preferences rational, continuous, convex, strictly convex, monotone, strongly monotone, locally nonsatiated?
\item Prove that $u$ is concave. Is $v$ concave? Does your answer change the shape of the preferences?
\item Is $w(x)=\mathbf 1\{u(x)\ge2\}$ another representation of the same preferences?
\end{enumerate}
\task{4. What exactly did we observe?}{Core · 15 minutes}
A consumer has a fixed weak ranking over $\{a,b,c\}$. In one observation they select $a$ from $\{a,b\}$; in another they select $b$ from $\{a,b,c\}$.
\begin{enumerate}[label=(\alph*)]
\item If each observation records only one selected maximiser, do these observations refute rational preferences? Construct a ranking if they do not.
\item Now the records are the \emph{full} best-alternative sets: $C(\{a,b\})=\{a\}$ and $C(\{a,b,c\})=\{b\}$. Can a rational ranking generate them? State the contradiction.
\item Which additional observation or recording convention would distinguish these interpretations?
\end{enumerate}
\task{5. Why does the budget line matter?}{Core · 15 minutes}
Let $B=\{x\in\R_+^2:p\cdot x\le m\}$, $p\gg0$, $m>0$.
\begin{enumerate}[label=(\alph*)]
\item Assuming a best bundle $x^*$ exists and preferences are LNS, prove $p\cdot x^*=m$ without derivatives.
\item Does budget exhaustion at every best bundle imply LNS? Try $u(x)=-[(x_1-2)^2+(x_2-2)^2]$ and $p=(1,1)$, $m=1$. Here ``every'' refers to every optimum for this one budget.
\item Does LNS alone guarantee that a best bundle exists? Try a lexicographic ranking that first maximises $x_1$ and, conditional on a tie, minimises $x_2$, on $B$ with $p=(1,1)$, $m=1$. Does \emph{this} example actually fail existence? If not, repair the example using $X=[0,1]$.
\end{enumerate}
\task{6. Can a finite grid repair lexicographic utility?}{Reserve · 10 minutes}
On $X_N=\{0,1/N,\ldots,1\}^2$, lexicographic preferences first compare $x_1$, then $x_2$. Find $M$ for which $u_N(x)=Mx_1+x_2$ represents them. Can the same finite $M$ work on all of $[0,1]^2$? Can \emph{any} real-valued utility represent lexicographic preferences on the square?
